% !TeX program = xelatex
% =====================================================================
% BAN TIENG VIET — HUJOS-TT (Tap chi Khoa hoc Dai hoc Hue: Ky thuat & Cong nghe)
% =====================================================================
% Cung mot noi dung va CUNG mot so lieu voi ban tieng Anh (main_hujos.tex):
%   * Bang so: tables_vi/ sinh TU DONG tu tables/ (ban EN, von sinh tu
%     outputs/*.csv) boi experiments/make_tables_hujos_vi.py — assert tap chu
%     so cua hai ban trung nhau, nhan dich qua bang tra cuong che.
%   * Than bai: blocks_vi/ dich tu paper2/blocks/ (ban EN 14 trang), giu nguyen
%     moi con so va moi hedge; thuat ngu theo glossary_vi.tex.
%   * Phu luc S1-S21: dung chung supplementary.tex (ban EN).
%
% Cau truc theo quy dinh HUJOS (muc 3): bai viet bang tieng Viet => tom tat +
% tu khoa TIENG VIET o DAU bai, ban TIENG ANH o CUOI bai. Day thu Tu theo dung
% bai tieng Viet da nop truoc do (hujos_submit.md, da qua audit_spec.py):
%   than bai -> DU LIEU, MA NGUON VA TUYEN BO -> TAI LIEU THAM KHAO ->
%   THONG TIN TAC GIA (+ chu ky) -> khoi tieng Anh.
%
% \HUJOSVI bat cong tac ngon ngu trong preamble.tex (caption "Bang/Hinh",
% dinh ly "Menh de/Ghi chu/Chung minh", header tieng Viet). Ban EN khong dinh
% nghia macro nay nen khong doi hanh vi.
% =====================================================================

\def\HUJOSVI{}
% !TeX program = xelatex
% =====================================================================
% HUJOS-TT (Hue University Journal of Science: Techniques and Technology)
% =====================================================================
% Moi thong so lay TU QUY DINH THAT cua tap chi (khong doan):
%   * https://jos.hueuni.edu.vn/index.php/hujos-tt/about/submissions
%       - toi da 12 trang, kho 19 cm x 27 cm
%       - font Palatino Linotype, 10 pt
%       - le 2 cm ca bon ben
%       - gian dong 1.15; 6 pt truoc va 3 pt sau moi doan
%       - cong thuc danh so ben phai, goi bang "Eq. (1)"
%       - caption BANG o TREN, caption HINH o DUOI, so bang/hinh IN DAM
%       - co khoang trong giua gia tri so va don vi (tru do/phut/giay)
%       - tai lieu tham khao kieu Vancouver (ngoac vuong) tu 2026
%       - abstract + tu khoa ca TIENG VIET va TIENG ANH, <= 250 tu
%       - tu khoa 3-5 cum, phai xuat hien trong abstract
%       - cuoi bai: ho ten, hoc ham, hoc vi, co quan, dia chi, dien thoai,
%         email, nguoi huong dan (neu co), chu ky tac gia, chi ro
%         corresponding author
%   * Template chinh thuc (HUJOS-NS, file DOCX tai tu Drive cua tap chi)
%     do tu word/document.xml va word/styles.xml:
%       - pgSz 11624 x 16840 twips = 20.5 x 29.7 cm, le tren/duoi 1.5 cm,
%         trai/phai 2 cm, HAI cot, columnsep 425 twips = 0.75 cm
%       - title 28 half-pt = 14 pt bold; tac gia 20 = 10 pt bold;
%         co quan 18 = 9 pt regular; abstract/tu khoa 18 = 9 pt
%       - Heading1 24 = 12 pt, before 360 twips = 18 pt, after 240 = 12 pt
%       - Heading2 before 18 pt, after 8 pt
%       - body spacing line 276 = 1.15, before 120 = 6 pt, after 60 = 3 pt
%       - font nhung trong template: PalatinoLinotype-{regular,bold,italic,
%         boldItalic}.ttf va Cardo
%     => Template NS khac quy dinh van ban cua TT/ARD ve kho giay va le.
%        Ta theo QUY DINH VAN BAN cua TT (19x27, le 2 deu) vi do la phan
%        bat buoc doi voi bai nop TT; giu HAI cot va columnsep 0.75 cm cua
%        template vi quy dinh van ban khong noi ve so cot.
%   * Kiem tra font: Palatino Linotype that (trich tu template) co 16/16
%     glyph tieng Viet theo cmap. `fc-list :lang=vi` bao KHONG ho tro tieng
%     Viet la SAI -- fontconfig doan theo bang lang chu khong doc cmap.
%     Dung TeX Gyre Pagella (clone Palatino, metric-compatible) de an toan
%     ban quyen font, da test compile tieng Viet 0 missing character.
% =====================================================================

\documentclass[10pt,twoside]{article}

% ---- hinh hoc ----
% MAU THUAN GIUA HAI NGUON CUA TAP CHI, da doi chieu ca hai:
%   (a) Quy dinh van ban o hujos-tt/about/submissions: "standard-sized paper
%       (19 cm x 27 cm) ... margins on all sides of 2 cm"  -> vung text 15 x 23 cm
%   (b) TEMPLATE DOCX chinh thuc (tai tu Drive tren trang authorguidelines), do tu
%       word/document.xml: <w:pgSz w:w="11624" w:h="16840"> = 20.5 x 29.7 cm (A4),
%       <w:pgMar w:top="851" w:bottom="851" w:left="1134" w:right="1134"> =
%       1.5 cm tren/duoi va 2 cm trai/phai, HAI cot, columnsep 425 twips = 0.75 cm
%       -> vung text 16.5 x 26.7 cm
% Chon (b) vi template la artifact ma toa soan thuc su dung de xep chu, va noi
% dung template ghi ro "format your manuscripts exactly the same as this document".
% Vung text (b) lon hon (a) 27.7% — do la cach hop phap nhat de dat tran 12 trang
% ma KHONG phai cat noi dung (cat noi dung la quyet dinh cua tac gia).
\usepackage[paperwidth=20.5cm,paperheight=29.7cm,
            top=1.5cm,bottom=1.5cm,left=2cm,right=2cm,
            headsep=0.4cm,footskip=0.6cm]{geometry}

% ---- font: Palatino (TeX Gyre Pagella) + math khop ----
\usepackage{fontspec}
\setmainfont{TeX Gyre Pagella}
\usepackage{unicode-math}
\setmathfont{TeX Gyre Pagella Math}

% ---- hai cot, columnsep 0.75 cm (425 twips theo template) ----
\setlength{\columnsep}{0.75cm}

% ---- gian dong 1.15, 6 pt truoc / 3 pt sau, khong gian giua doan = 0 ----
\linespread{1.15}
\setlength{\parskip}{0pt}
\setlength{\parindent}{0.45cm}

% ---- toan hoc ----
% KHONG dung \numberwithin{equation}{section}: quy dinh HUJOS-TT yeu cau cong thuc
% danh so lien tuc va duoc goi la "Eq. (1)", nen danh so theo section se ra (1.1)
% va lam sai ca cach goi lan thu tu trong van ban.
%
% KHONG nap amssymb CUNG unicode-math. Da do bang build that: hai goi trung dinh
% nghia 4 lenh va dung bien dich voi 4 loi fatal
%   ! LaTeX Error: Command `\eth' already defined.
%   ! LaTeX Error: Command `\smallsetminus' already defined.
%   ! LaTeX Error: Command `\digamma' already defined.
%   ! LaTeX Error: Command `\backepsilon' already defined.
% unicode-math (nap o tren, kem TeX Gyre Pagella Math) da cung cap day du symbol
% cua amssymb, nen bo amssymb la cach sua dung thay vi \let ghi de tung lenh.
\usepackage{amsmath}

% ---- bang: booktabs + caption tren bang ----
\usepackage{booktabs}
\usepackage[font=small,labelfont=bf,labelsep=period,
            skip=4pt,justification=raggedright,singlelinecheck=false]{caption}
\captionsetup[table]{position=top}
\captionsetup[figure]{position=bottom}

% ---- heading theo so do tu template ----
\usepackage{titlesec}
\titleformat{\section}{\normalfont\bfseries\fontsize{12}{14}\selectfont}
  {\thesection.}{0.5em}{}
\titlespacing*{\section}{0pt}{18pt}{12pt}
\titleformat{\subsection}{\normalfont\bfseries\fontsize{10}{12}\selectfont}
  {\thesubsection.}{0.5em}{}
\titlespacing*{\subsection}{0pt}{18pt}{8pt}
\titleformat{\subsubsection}[runin]{\normalfont\itshape}
  {\thesubsubsection.}{0.5em}{}[.]
\titlespacing*{\subsubsection}{0pt}{12pt}{0.5em}
% 18 pt space after text in each section (template yeu cau)
\titleformat{name=\section,numberless}{\normalfont\bfseries\fontsize{12}{14}\selectfont}
  {}{0pt}{}
\titlespacing*{name=\section,numberless}{0pt}{18pt}{12pt}

% ---- dinh ly / chung minh ----
\usepackage{amsthm}
\newtheorem{proposition}{Proposition}
\newtheorem{remark}{Remark}
\renewcommand{\proofname}{Proof}

% ---- trích dẫn Vancouver (ngoặc vuông) từ 2026 ----
\usepackage[numbers,sort&compress]{natbib}
\bibpunct{[}{]}{,}{n}{}{,}

% ---- link (an vien, khong doi mau) ----
% PHAI \PassOptionsToPackage TRUOC hyperref: hyperref tu nap url ben trong, nen
% viet \usepackage[hyphens]{url} o SAU se dung "Option clash for package url"
% va dung bien dich. Da tra gia: build bao P3=1 ma van xuat PDF 13 trang, trong
% do 3 overfull cua .bbl giu nguyen 189.81/139.74/118.98pt -> hyphens chua an.
\PassOptionsToPackage{hyphens}{url}
\usepackage[hidelinks]{hyperref}

% ---- header/footer nhu template: khong tu them so trang ngoai cac thu da set ----
\usepackage{fancyhdr}
\pagestyle{fancy}
\fancyhf{}
\renewcommand{\headrulewidth}{0pt}
\fancyhead[LE]{\footnotesize\itshape X. V. Mai and T. T. Nguyen}
\fancyhead[RO]{\footnotesize\itshape Hue University Journal of Science: Techniques and Technology}
\fancyfoot[C]{\footnotesize\thepage}

% ---- phu luc: doi ten "Appendix" thanh "Appendix" (giu nguyen) ----
\usepackage{appendix}

% ---- tien ich ----
% [hyphens] la BAT BUOC: ba overfull nang nhat cua bai (189.81 / 139.74 /
% 118.98 pt, tuc chu loi ra ngoai trang) deu nam trong .bbl va deu do URL dai
% khong ngat duoc. Mac dinh goi url chi ngat o dau '/' va '.', ma ten mien nhu
% 'technical-explainer-carbon-credits-for-improved-rice-cultivation' khong co '/'.
% Option duoc truyen bang \PassOptionsToPackage o tren (truoc hyperref) de tranh
% option clash, nen o day nap khong kem option.
\usepackage{url}
\usepackage{xspace}
\usepackage{microtype}
\usepackage{array}

% graphicx la BAT BUOC cho \resizebox. Thieu no thi 24 loi fatal
% "Undefined control sequence" + "Missing number" + "Illegal unit of measure"
% (vi \resizebox bi phan tich sai). Da tra gia: ban IEEE co goi nay, ban HUJOS
% quen -> build ra PDF 13 trang nhung la RAC, khong duoc bao cao la thanh cong.
\usepackage{graphicx}

% \fittocol: chi THU NHO bang khi no that su rong hon cot, KHONG BAO GIO phong to.
%
% Tai sao khong dung \sbox de do: cot p{} chua \par, ma \sbox = \savebox trong
% \hbox nen khong do duoc vat lieu dang doan van. Do bang \sbox cho ra 7.92-8.01cm
% trong khi build that bao overfull 18-26pt (~8.8cm) -> so do sai, dan den ket luan
% sai rang "bang khong tran". lrbox la cong cu dung cho hop chua \par.
%
% Tai sao khong dung \resizebox{\columnwidth}{!} truc tiep: no CUONG BUC dung bang
% kho cot, nen bang hep hon (vd tab_baseline 13.69cm trong table* 16.5cm) se bi
% PHONG TO ~20%, chu to hon thiet ke.
\newsavebox{\fitbox}
\newcommand{\fittocol}[1]{%
  \begin{lrbox}{\fitbox}#1\end{lrbox}%
  \ifdim\wd\fitbox>\columnwidth
    \resizebox{\columnwidth}{!}{\usebox{\fitbox}}%
  \else
    \usebox{\fitbox}%
  \fi}

% Cho phep keo gian khoang trang them 2em khi khong tim duoc diem ngat.
% Ly do: cac doan chua nhieu \ref lien tiep (vd
% "(Sec.~\ref{sec:umax}, Sec.~\ref{sec:v2reach}, Table~\ref{tab:umax},
% Table~\ref{tab:f2reach})") tao ra chuoi khong ngat duoc, gay overfull
% 34.89pt va 26.16pt trong 01_intro.tex. Day la cach sua khong doi noi dung.
\setlength{\emergencystretch}{2em}

% Kho cot that voi geometry nay: text 15 cm, 2 cot, columnsep 0.75 cm
%   -> moi cot = (16.5 - 0.75)/2 = 7.875 cm
%   -> table* (2 cot) = 16.5 cm
% Day la con so ma experiments/make_tables_hujos.py dung de gate chieu rong.
\newlength{\HUJOScol}
\setlength{\HUJOScol}{7.875cm}

% Front matter: title 14pt bold, authors 10pt bold, affiliations 9pt regular,
% abstract + keywords 9pt (dung so do tu template)
\newcommand{\hujostitle}[1]{%
  {\centering\normalfont\bfseries\fontsize{14}{17}\selectfont #1\par}}
\newcommand{\hujosauthors}[1]{%
  {\raggedright\normalfont\bfseries\fontsize{10}{12}\selectfont #1\par}}
\newcommand{\hujosaffil}[1]{%
  {\raggedright\normalfont\fontsize{9}{11}\selectfont #1\par}}
\newenvironment{hujosabstract}
  {\par\vspace{10pt}\normalfont\fontsize{9}{11}\selectfont\noindent\ignorespaces}
  {\par}


\begin{document}

% ---------------------------------------------------------------- front matter VI
\begingroup
  \hujostitle{Nhật ký tưới được neo theo thiết bị: đánh giá đối kháng phương án
  MRV chi phí gần như bằng không cho lúa phát thải thấp}

  \vspace{8pt}
  \hujosauthors{Mai Xuân Văn$^{1}$, Nguyễn Tương Tri$^{2,3,*}$}

  \vspace{4pt}
  \hujosaffil{$^{1}$ Trường Trung học phổ thông (THPT) Thuận Hóa, Trường Đại học Sư phạm, Đại học Huế, Huế, Việt Nam\\
  $^{2}$ Viện Đào tạo Mở và Công nghệ thông tin, Đại học Huế, Huế, Việt Nam\\
  $^{3}$ Trường Đại học Sư phạm, Đại học Huế, 34 Lê Lợi, Huế, Việt Nam}

  \vspace{4pt}
  \hujosaffil{$^{*}$ Tác giả liên hệ: Nguyễn Tương Tri. E-mail: ntuongtri@hueuni.edu.vn}

  \vspace{6pt}
  \begin{hujosabstract}
  % TRAN 250 TU: ban dau 316 tu (vuot 26%). Da cat nhung cho lap y voi than bai
  % (danh sach 6 truc vi pham mo hinh, giai thich "do chinh la dong gop"), giu
  % NGUYEN VEN moi con so va moi hedge. Dem bang script check_abstract_words.
  \noindent\textbf{Tóm tắt.}

  Đo đạc, báo cáo, thẩm định (MRV) là nút thắt khiến mức giảm phát thải của lúa phát thải thấp chưa thành tín
  chỉ các-bon; giám sát chiếm ba phần tư chi phí. Vì mọi hành động quản lý nước
  đủ điều kiện cấp tín chỉ trong tưới khô ướt luân phiên (AWD) vốn là lệnh điều
  khiển, nhật ký nông dân tự ghi trông như bằng chứng gần như miễn phí, xếp hạng cao
  nhất theo Quyết định 4801/QĐ-BNNMT. Chúng tôi dựng lớp đó trên mô hình sinh đôi số
  FAO-56 rồi tấn công nó; phần tấn công là đóng góp. Phát lại nhật ký với
  ERA5 công khai phơi bày độ sâu không đất nào chứa được, nhưng đối thủ giữ
  nguyên tổng lượng và từng độ sâu, chỉ dịch thời điểm lệnh, làm số pha khô được cấp tín
  chỉ tăng $33\%$ (ghép cặp). Không lớp vật lý nào thoát hữu ích: bộ phát hiện
  mô hình điểm kết tội oan $26{,}7\%$ nông dân trung thực, còn bộ bền vững giảm oan xuống
  $5{,}0\%$ lại bỏ lọt phần lớn ca dịch thời điểm. Trên $281$ nhật ký tại ba trạm
  Đồng bằng sông Cửu Long, hai năm khí hậu, đối chiếu với bản ghi lệnh có chữ ký, dấu thời gian của gateway bắt $116/116$ nhật ký giả, không oan ai. Vật lý không thể
  đứng một mình; dấu thời gian phải do thiết bị cấp, không phải nông dân.

\noindent\textbf{Từ khóa:} tưới khô ướt luân phiên; đo đạc, báo cáo, thẩm định; tín chỉ
  các-bon; lúa phát thải thấp; mô hình sinh đôi số.

  \vspace{10pt}
  \end{hujosabstract}
\endgroup

% ---------------------------------------------------------------- than bai
% BAN TIENG VIET cua 01_intro.tex — dich tu ban EN 14 trang (2026-10-08).
% Moi con so giu nguyen, chi doi sang dinh dang Viet Nam (dau phay thap phan).
% Thuat ngu theo glossary_vi.tex. KHONG them/bot claim so voi ban EN.
\section{Mở đầu}
\label{sec:intro}

Việt Nam đã cam kết một triệu héc-ta lúa phát thải thấp vào năm 2030, và
phần nông học đã ngã ngũ: tăng độ khắc nghiệt của phơi hạn ngắt quãng làm giảm
phát thải khí mê-tan từ $41$--$71\%$ so với ngập liên tục mà không làm giảm
năng suất hạt~\cite{balaine2019}, và chế độ nước là nhân tố chi phối thông
lượng mê-tan theo mùa trên các ruộng lúa châu Á~\cite{yan2005}. Điều chưa ngã
ngũ là đo lường. Theo Quyết định 4801/QĐ-BNNMT, bằng chứng được xếp hạng cao
nhất cho một yêu cầu cấp tín chỉ là bản ghi các hoạt động tưới và rút nước
đã thực hiện~\cite{qd4801}, và theo VM0051 của Verra, chính bản ghi đó phải
vượt qua thẩm định của bên thứ ba~\cite{vm0051}. Làm ra nó tốn khoản tiền mà
một hợp tác xã nông hộ nhỏ không có: riêng giám sát chiếm khoảng $77\%$ chi
phí MRV~\cite{jem2026mrv}, và quy trình hiện hành dựa vào việc chụp ảnh định
kỳ một ống đo mực nước.

Cách sửa hiển nhiên là nhận ra rằng một lệnh tưới hiện đại vốn đã là một sự
kiện số: nếu gateway đóng mở máy bơm thì nó cũng có thể ghi lại lần đóng mở
đó, và bản ghi ấy là dữ liệu hoạt động bậc 1 (Tier-1) theo nghĩa của Quyết
định 4801 với chi phí biên bằng không --- đây là tiền đề của công trình đồng
hành~\cite{twingate2026}, vốn dùng một mô hình sinh đôi số FAO-56 làm cổng an
toàn quyết định một lệnh có được truyền qua kênh có tổn thất hay không. Chúng
tôi đi thêm một bước và đặt câu hỏi: liệu chính mô hình sinh đôi số đó có thể
\emph{kiểm toán} những gì nông dân khai báo sau vụ việc hay không, mà không
cần triển khai cảm biến nào: một nhật ký các độ sâu khai báo được phát lại
qua cân bằng nước dưới dữ liệu ERA5 công khai, và một khai báo mà đất không
thể nào hấp thụ được là bằng chứng của sự ngụy tạo. Phiên bản đầu tiên của
lớp này báo cáo không kết tội oan nhật ký trung thực nào trong $60$ nhật ký
và phát hiện $60/60$ nhật ký giả --- đúng nhưng vô giá trị, vì các quỹ đạo
được dán nhãn ``thật'' lại do chính mô hình sau đó kiểm toán chúng sinh ra:
bộ phát hiện đang tự chấm bài của mình (Phụ lục~S10.4).

Bài báo này trình bày điều xảy ra khi vòng luẩn quẩn đó bị gỡ bỏ: chúng tôi
dựng lại ruộng sao cho nó vi phạm mọi giả định của mô hình sinh đôi số, và
thay các kiểu gian lận ngây thơ bằng một đối thủ đã đọc bộ phát hiện. Ba phát
hiện sau đây, đồng thời là ba đóng góp của bài. Thứ nhất, bộ phát hiện mô
hình điểm không chỉ lạc quan, nó \emph{không an toàn}: nó kết tội oan
$26{,}7\%$ nông dân trung thực, vì một ngưỡng tuyệt đối trên chênh lệch bơm
đối xử với người quên một mục ghi y hệt người tưới trộm; quy tắc bền vững sửa
được lỗi này --- một phép giao trên tập bất định $\Theta$ gồm $30$ cấu hình
đất, hệ số cây trồng và hiệu suất (Mục~\ref{sec:robust}) --- phải trả giá
bằng năng lực phát hiện, điều chúng tôi đo lường chứ không che giấu. Thứ
hai, việc chứng thực dựa trên vật lý có một lỗ hổng mà không lượng mô hình
hoá nào bịt được: một đối thủ giữ nguyên tổng lượng nước và từng độ sâu riêng
lẻ có thể \emph{dịch lệnh theo thời gian}, làm thổi phồng các pha khô đủ
điều kiện cấp tín chỉ --- đúng đại lượng mà VM0051 và Quyết định 4801 trả
tiền --- thêm $+33\%$ trong so sánh ghép cặp ($t=5{,}29$; $+38\%$,
$t=3{,}46$ với tập bất định đầy đủ). Không tầng vật lý nào thoát được --- mô
hình điểm bắt được bốn trong chín nhật ký bị sửa nhưng lại chính là bộ phát
hiện kết tội oan $26{,}7\%$ nông dân trung thực, trong khi tầng bền vững hạ
tỷ lệ oan xuống $5{,}0\%$ thì để lọt tám trong chín --- và chúng tôi chứng
minh điều này không thể tránh khỏi với mọi bộ phát hiện là hàm của quỹ đạo
phát lại (Mục~\ref{sec:threat}, Mục~\ref{sec:timeshift}). Thứ ba, cách sửa
rẻ và không mang tính vật lý: đối chiếu nhật ký với bản ghi lệnh có chữ ký và
dấu thời gian của gateway bắt được $116/116$ nhật ký thật sự bị giả mạo và
không kết tội oan nhật ký nào trong $60$ nhật ký trung thực, với chi phí
không vượt quá chính gateway mà công trình đồng hành đã triển khai
(Mục~\ref{sec:anchor}). Bài báo cũng sửa lại ba con số đã công bố trước đó:
độ sâu chấp nhận được là một dải $[49{,}0; 86{,}0]$~mm chứ không phải một
giá trị $61$~mm; bộ lọc vật lý thứ ba không thể kích hoạt trên lúa vụ khô
đồng bằng; và ngưỡng tuyệt đối trên chênh lệch bơm là không thể chấp nhận
(Mục~\ref{sec:umax}, Mục~\ref{sec:v2reach}).

% BAN TIENG VIET cua 07_system_model.tex (2026-10-08). So lieu giu nguyen.
\section{Mô hình hệ thống}
\label{sec:model}

\subsection{Cân bằng nước của ruộng}

Lượng nước trong tầng rễ $w_t$ [mm] biến đổi theo giờ dưới mô hình thùng chứa
FAO-56~\cite{fao56}:
\begin{equation}
w_{t+1}=\mathrm{clip}_{\le FC}\!\left(w_t+P_t-RO_t-ET_{c,t}+u_t\right),
\label{eq:bucket}
\end{equation}
với lượng bốc thoát hơi nước của cây trồng và hệ số ứng suất nước
\begin{equation}
ET_{c,t}=K_{s,t}\,K_c\,ET_{0,t},
\qquad
K_{s,t}=\frac{FC-w_t}{(1-p)\,\mathrm{TAW}},
\label{eq:etc}
\end{equation}
trong đó $ET_{0,t}$ là bốc thoát hơi nước tham chiếu theo giờ
(Penman-Monteith), $P_t$ là lượng mưa, $RO_t$ là dòng chảy mặt (phương pháp
đường cong số SCS, tính theo ngày rồi chia đều cho các giờ của một trận mưa),
và $u_t$ là độ sâu tưới áp dụng trong giờ $t$. Lượng nước vượt quá sức chứa ẩm
đồng ruộng $FC$ bị toán tử cắt bỏ loại đi, nên bảo toàn khối lượng đúng tuyệt
đối và được kiểm tra như một bất biến. Độ cạn kiệt của tầng rễ là
\begin{equation}
d_t=\frac{FC-w_t}{FC-WP},\qquad \mathrm{span}=FC-WP,
\label{eq:depl}
\end{equation}
với $WP$ là điểm héo; đối với loại đất thịt pha sét limon tham chiếu của đồng
bằng, $\mathrm{TAW}=FC-WP=70$~mm, $FC=160$~mm, $WP=90$~mm.

\subsection{Phát lại một nhật ký khai báo}
\label{sec:replay}

Giả sử một nhật ký khai báo cung cấp số hạng tưới của \eqref{eq:bucket}.
\emph{Toán tử phát lại} $\mathcal{R}$ ánh xạ một nhật ký sang quỹ đạo mà nó
hàm ý, bằng cách lặp \eqref{eq:bucket} dưới dữ liệu khí tượng quan trắc
$(P_t,ET_{0,t})$:
\begin{equation}
\mathcal{R}(L)=\{(\hat w_t,\hat d_t)\}_{t=0}^{T},
\label{eq:replay}
\end{equation}
với $\hat w_0=w_{\mathrm{init}}$ và
\begin{equation}
u_t=\sum_{(t_i,u_i)\in L}u_i\,\mathbb{1}[t_i=t],
\label{eq:uterm}
\end{equation}
sao cho $\hat w_t$ là lượng nước tầng rễ mà nhật ký \emph{khai} rằng ruộng đã
có, và $\hat d_t$ là độ cạn kiệt của nó theo \eqref{eq:depl}; dấu mũ đánh dấu
các đại lượng được tái dựng từ nhật ký chứ không đo được. Việc phát lại là
tất định và chỉ dùng dữ liệu tái phân tích công khai, nên hai kiểm toán viên
phát lại cùng một nhật ký sẽ thu được cùng một quỹ đạo và phán quyết có thể
tái lập. Dấu mũ trừ chỉ trạng thái ngay \emph{trước} khi khai báo của giờ đó
được áp dụng, như trong $\hat w_{t}^{-}$: lượng nước sẵn có để tiếp nhận một
lần tưới, và đó là thứ bộ lọc phải xem xét.

Toán tử này cũng là thứ làm cho phương án trở nên rẻ: không cảm biến nào bị
truy vấn, và đầu vào duy nhất là nhật ký, chuỗi thời tiết công khai cùng một
hồ sơ đất. Điểm yếu của nó, được trình bày ở Mục~\ref{sec:timeshift}, là nó
chỉ kiểm tra được một quỹ đạo có \emph{khả dĩ} hay không, chứ không bao giờ
biết được nó có \emph{đã xảy ra} hay không.

% BAN TIENG VIET cua 02_threat_attack.tex (2026-10-08). So lieu giu nguyen.
\section{Mô hình đe doạ}
\label{sec:threat}

Một nhật ký tưới khai báo là một ánh xạ hữu hạn
\begin{equation}
L=\{(t_i,u_i)\}_{i=1}^{|L|},\qquad u_i>0,\ t_i\in[0,T),
\label{eq:log}
\end{equation}
trong đó $u_i$ là độ sâu khai báo tính bằng milimét và $t_i$ là chỉ số giờ.
Ba đại lượng độc lập có thể bị làm giả, và mỗi tầng phòng thủ quan sát một tập
con khác nhau của chúng:
\begin{equation}
\underbrace{\textstyle\sum_i u_i}_{\text{tổng lượng}}\,,\qquad
\underbrace{\{u_i\}}_{\text{từng độ sâu}}\,,\qquad
\underbrace{\{t_i\}}_{\text{dấu thời gian}}\,.
\label{eq:three}
\end{equation}
Bảng~S1 của Phụ lục liệt kê các lớp gian lận tương ứng. Đối thủ là
\emph{duy lý}: nó tối đa hoá số pha khô được cấp tín chỉ với điều kiện không
bị đánh dấu, và nó biết bộ phát hiện, vì bộ phát hiện đã được công bố; đó là
nghĩa mà theo đó $\varepsilon$ trong \eqref{eq:v4} bên dưới không phải là bí
mật.

Gọi $M$ là lượng nước thật sự đã bơm, đo tại trạm của hợp tác xã, và
$U(L)=\sum_i u_i$ là tổng lượng khai báo; thống kê của tầng~1 là
\begin{equation}
G=M-U(L),
\label{eq:gap}
\end{equation}
và đại lượng liên quan đến tín chỉ thì đếm \emph{số pha khô}. Viết
$\mathcal{A}(\cdot)$ cho tập các khoảng tối đại $[a,b)$ của một quỹ đạo phát
lại mà trên đó độ cạn kiệt luôn ở mức từ $d_{\mathrm{dry}}$ trở lên và không
có khai báo tưới nào, kéo dài ít nhất $72$~giờ --- mức tối thiểu ba ngày của
Quyết định 4801~\cite{qd4801}:
\begin{equation}
\begin{aligned}
\mathcal{A}\bigl(\{\hat d_t\}\bigr)=\Bigl\{[a,b):\ &\min_{t\in[a,b)}
\hat d_t\ge d_{\mathrm{dry}},\\
&u_t=0\ \forall t\in[a,b),\\
&b-a\ge 72\,\mathrm{h}\Bigr\}.
\end{aligned}
\label{eq:awd}
\end{equation}
Số pha được cấp tín chỉ khi đó là
\begin{equation}
n_{\mathrm{dry}}(L)=\bigl|\mathcal{A}\bigl(\mathcal{R}(L)\bigr)\bigr|,
\label{eq:ndry}
\end{equation}
với $d_{\mathrm{dry}}=0{,}45$ trong toàn bộ bài. Doanh thu tỉ lệ với
$n_{\mathrm{dry}}$ chứ không tỉ lệ với $U(L)$ --- chính sự phân biệt này làm
cho tấn công ở Mục~\ref{sec:timeshift} sinh lợi, và nó vô hình trong một
cách phát biểu coi tổng lượng nước là mục tiêu.

\section{Các bộ lọc dựa trên vật lý và đúng điểm mù của chúng}
\label{sec:filters}

Phát lại $L$ dưới dữ liệu ERA5 công khai cho ra $(\hat w_t,\hat d_t)$ (công
trình đồng hành~\cite{twingate2026} trình bày đầy đủ cân bằng này); có ba bộ
lọc được áp dụng. \emph{$V_1$, bão hoà}: một khai báo phát ra khi tầng rễ đã
ở sức chứa ẩm đồng ruộng thì không thể được lưu trữ:
\begin{equation}
V_1(L)=\{\,i:\ u_i>0\ \wedge\ \hat w_{t_i}^{-}\ \ge\ FC-\varepsilon\,\}.
\label{eq:v1}
\end{equation}

\emph{$V_4$, vượt sức chứa}: một khai báo lớn hơn phần thiếu hụt sẵn có tại
giờ đó hẳn đã chảy tràn hoặc thấm mất, nên nó không ``được áp dụng'' theo
nghĩa mà quy chuẩn đòi hỏi:
\begin{equation}
V_4(L)=\{\,i:\ u_i\ >\ \underbrace{(FC-\hat w_{t_i}^{-})}_{\text{phần thiếu hụt}}
\ +\ \varepsilon\,\}.
\label{eq:v4}
\end{equation}

\emph{$V_5$, giới hạn thấm.} Không phụ thuộc trạng thái, một lần tưới kéo dài
$\tau_i$ không thể cấp nhiều hơn lượng đất tiếp nhận được:
\begin{equation}
u_i\ \le\ k_{\mathrm{inf}}\,\tau_i+p_{\max},
\label{eq:v5}
\end{equation}
với $k_{\mathrm{inf}}=10$~mm/h và $p_{\max}=100$~mm, lấy ở đầu rộng của dải
giá trị đã công bố cho ruộng lúa làm dầm (giới hạn $110$~mm tại
$\tau_i=1$~giờ), để một nông dân ghi hai lần tưới liền kề trong cùng một giờ
không bị kết tội.

Bộ lọc thứ tư, về độ cạn kiệt xuống dưới điểm héo, đã có trong phiên bản công
bố của lớp này. Nó không thể kích hoạt: xem Mục~\ref{sec:v2reach}.

\begin{proposition}[Điểm mù chính xác của các phòng thủ tổng hợp]
\label{prop:blind}
Giả sử $L$ và $L'$ có cùng đa tập hợp các độ sâu và cùng tổng lượng, tức
$\{u_i\}=\{u'_j\}$ theo nghĩa đa tập hợp. Khi đó
\begin{equation}
G(L)=G(L')\qquad\text{và}\qquad |V_5(L)|=|V_5(L')|.
\end{equation}
\end{proposition}

\begin{proof}
Xem Phụ lục~S11; cả hai đại lượng đều là hàm đối xứng của đa tập hợp độ sâu
và không bao giờ phụ thuộc $t_i$.
\end{proof}

Mệnh đề~\ref{prop:blind} gói trọn toàn bộ khó khăn trong một dòng: hai phòng
thủ không phụ thuộc việc mô hình có đúng hay không --- đồng hồ đo bơm và giới
hạn thấm --- đều bất biến \emph{chính xác} dưới phép biến đổi mà đối thủ cần.
Còn lại $V_1$ và $V_4$, vốn phụ thuộc thứ tự qua $\hat w_{t_i}^{-}$; hoá ra
chúng không giúp được gì.

\section{Tấn công dịch thời điểm}
\label{sec:timeshift}

\begin{proposition}[Lợi ích tín chỉ và sự lọt lưới là cùng một hành động]
\label{prop:coincide}
Giả sử một khai báo tại giờ $t$ có phần thiếu hụt sẵn có
$c_t=FC-\hat w_t^{-}$. Giữa hai khai báo liên tiếp, thùng chứa mất đi
$\Delta_t=\sum_{s}(\mathit{ETc}_s+\mathit{DP}_s-P_s^{+})\ge 0$, nên việc hoãn
một khai báo thêm $\delta$ giờ làm tăng phần thiếu hụt của nó đúng bằng lượng
nước đã thoát đi trong các giờ đó:
\begin{equation}
c_{t+\delta}\ \ge\ c_t+\min\!\bigl(\delta\cdot\overline{\Delta},\ FC-\hat w_t^{-}\bigr).
\label{eq:delay}
\end{equation}
Do đó phép biến đổi dời các khai báo \emph{ra khỏi} một khoảng khô và đặt lên
biên của khoảng đó đồng thời (i)~làm tăng $n_{\mathrm{dry}}$, vì khoảng đó
không còn bị ngắt quãng, và (ii)~làm tăng $c_t$ tại mỗi khai báo bị dời, điều
mà theo \eqref{eq:v4} khiến $V_4$ \emph{ít} có khả năng kích hoạt hơn. Đối
thủ không phải đánh đổi điều gì: gradient của lợi nhuận và gradient của sự
không thể phát hiện cùng chỉ một hướng.
\end{proposition}

\begin{proof}
Xem Phụ lục~S11. Nói ngắn gọn: (i) đúng theo định nghĩa, vì $\mathcal{A}$ đòi
hỏi không có khai báo nào bên trong khoảng; (ii) suy ra được vì hoãn một khai
báo để thùng chứa tiếp tục thoát nước thêm $\delta$ giờ mà không có dòng vào,
nên phần thiếu hụt sẵn có chỉ có thể tăng trong khi $u_i$ không đổi.
\end{proof}

Đó là lý do tấn công này không thể bị bịt bằng thuỷ văn tốt hơn: bất kỳ bộ
phát hiện nào có phán quyết là hàm của quỹ đạo phát lại đều phải tích phân
đúng lượng nước thoát đã làm cho khai báo bị hoãn trông có vẻ hợp pháp. Bản
thực nghiệm của tấn công là một thuật toán leo đồi trên các độ lệch giờ
nguyên trong $[-36,36]$, chỉ chấp nhận một bước dời nếu tổng lượng được bảo
toàn từng bit, mọi độ sâu vẫn ở dưới \eqref{eq:v5}, và $n_{\mathrm{dry}}$
tăng nghiêm ngặt. Tấn công không ăn cắp nước --- thứ bị ăn cắp là
\emph{lịch sử} (Phụ lục~S19).


% BAN TIENG VIET cua 03_guarantee.tex (2026-10-08). So lieu giu nguyen.
\section{Khi nào được kết tội một nông dân?}
\label{sec:guarantee}

\subsection{Độ sâu chấp nhận được là một dải, không phải một con số}
\label{sec:umax}

Phiên bản đã công bố của lớp này nêu một giới hạn duy nhất, $u\le 61$~mm, là
khai báo lớn nhất không bao giờ gây báo động --- đó là sản phẩm của một hồ sơ
đất và một điều kiện đầu duy nhất. Với $V_4$ kích hoạt khi
$u_i>c_{t_i}+\varepsilon$ và $c_t=d_t\cdot\mathrm{span}$, giới hạn tại một
ruộng có độ cạn kiệt ban đầu $w_0$ và tưới không muộn hơn $d^{*}$ là
\begin{equation}
u_{\max}=\max(d^{*},w_0)\cdot\mathrm{span}+\varepsilon .
\label{eq:umax}
\end{equation}
Hai điều chỉnh rút ra, cả hai đều được xác minh bằng cách quét toán tử phát lại
chứ không chỉ bằng đại số (dẫn xuất ở Phụ lục~S14): \eqref{eq:umax} dùng
$\max$ chứ không dùng hiệu --- một ruộng khô ban đầu chấp nhận khai báo
\emph{lớn hơn} ruộng ướt, và biểu thức đã công bố trước đó không tương ứng với
trạng thái vật lý nào --- và nó chặn một \emph{lần tưới đơn lẻ} chứ không chặn
cả vụ, nên một đối thủ phát nhiều khai báo dưới ngưỡng có thể thổi phồng tổng
lượng mà không có giới hạn. Bảng~S2 của Phụ lục đưa ra dải phải được công bố
thay thế, $[49{,}0; 86{,}0]$~mm trên toàn $\Theta$.

\subsection{Kết tội bền vững dưới bất định mô hình}
\label{sec:robust}

Kiểm toán viên không biết đất của nông dân, không biết giai đoạn sinh trưởng,
cũng không biết hiệu suất của hệ thống dẫn nước. Gọi $\Theta$ là tập bất định
và $V_\theta(L)$ là tập vi phạm dưới cấu hình $\theta$; bộ phát hiện đã công bố
chỉ đánh giá một $\theta_0$ mặc định và coi mọi vi phạm là có tội. Chúng tôi
thay nó bằng một phép giao:
\begin{equation}
\begin{aligned}
V_{\mathrm{rob}}(L)&=\bigcap_{\theta\in\Theta}V_\theta(L),\\
\Theta&=\mathcal{S}\times\mathcal{K}\times\mathcal{E}\times\mathcal{W},
\end{aligned}
\label{eq:robust}
\end{equation}
với $|\mathcal{S}|=5$ loại đất, $|\mathcal{K}|=2$ chế độ hệ số cây trồng,
$|\mathcal{E}|=3$ mức hiệu suất tưới và $|\mathcal{W}|=5$ trạng thái đầu, nên
$|\Theta|=150$ và mỗi nhật ký tốn $150$ lần phát lại. Ngữ nghĩa được chọn cố ý
lệch về một phía:
\begin{equation}
a(L)=
\begin{cases}
\text{kết tội}, & V_{\mathrm{rob}}(L)\neq\emptyset,\\
\text{chờ xác minh}, & \text{trường hợp còn lại}.
\end{cases}
\label{eq:verdict}
\end{equation}

\begin{proposition}[Không kết tội oan trong $\Theta$]
\label{prop:noalarm}
Nếu cấu hình thật của ruộng $\theta^{*}$ thuộc $\Theta$ và $L$ tuân thủ dưới
$\theta^{*}$, thì $V_{\mathrm{rob}}(L)=\emptyset$ và nông dân không bị kết tội.
\end{proposition}

\begin{proof}
$V_{\mathrm{rob}}(L)\subseteq V_{\theta^{*}}(L)=\emptyset$.
\end{proof}

Mệnh đề~\ref{prop:noalarm} là thứ làm cho lớp này triển khai được: sai số mô
hình có thể làm giảm năng lực phát hiện nhưng không thể tạo ra một lời kết
tội, với một cái giá phải được nói rõ. Hai chi tiết cài đặt --- cách xử lý các
thế giới không kích hoạt và độ trải của $\mathcal{W}$ --- sẽ âm thầm phá vỡ
mệnh đề nếu làm sai (Phụ lục~S7).

\subsection{Ngưỡng chênh lệch bơm tỉ lệ thuận}

Quy tắc tuyệt đối $G>\tau_0$ với $\tau_0=5$~mm giả định nhật ký là bản sao
trung thực của lượng nước đã bơm. Nó không phải vậy: quên ghi, làm tròn đến
$5$~mm gần nhất và ghi trễ một ngày đều đóng góp vào $G$ mà không có hành vi
sai trái nào, nên ngưỡng phải co giãn theo số lần tưới được ghi:
\begin{equation}
\begin{aligned}
\tau(n)&=\delta_{\mathrm{rel}}M+z\sqrt{n}\,\sigma_{\mathrm{event}},\\
\sigma_{\mathrm{event}}^{2}&=p(1-p)\bar u^{2}+\tfrac{r^{2}}{12}
+q(1-q)(s\bar u)^{2},
\end{aligned}
\label{eq:tau}
\end{equation}
với $p$ là tỷ lệ quên, $r$ là bước làm tròn, $q$ là tỷ lệ sai độ sâu và $s$ là
biên độ của nó, $\bar u=30$~mm là độ sâu điển hình của một lần tưới,
$\delta_{\mathrm{rel}}=2\%$ và $z=3$; quy tắc tuyệt đối là trường hợp suy biến
$n\to0$, $\sigma_{\mathrm{event}}\to0$.

\begin{remark}[Một bộ lọc không thể kích hoạt]
\label{sec:v2reach}
Bộ lọc vật lý thứ ba của phiên bản đã công bố, đánh dấu độ cạn kiệt xuống dưới
điểm héo, đã không hoạt động theo hai cách: ngưỡng của nó ($1{,}02$) nằm ngoài
dải đã cắt ngọn của chỉ số được báo cáo ($0$ lần kích hoạt trong $281$ nhật
ký), và sau khi sửa thì nó vẫn không thể kích hoạt --- qua $150$ cấu hình với
một nhật ký rỗng, độ cạn kiệt chưa cắt ngọn chỉ đạt đỉnh $0{,}9987$, và một
ruộng lúa vụ khô ở đồng bằng không chạm tới điểm héo (Bảng~S3 của Phụ lục, dẫn
xuất ở S13). Lớp này có hai bộ lọc vật lý hoạt động và một thiết bị an toàn
bất động.
\end{remark}


% BAN TIENG VIET cua 04_anchor.tex (2026-10-08). So lieu giu nguyen.
\section{Chứng thực neo theo thiết bị}
\label{sec:anchor}

\subsection{Bản ghi của gateway}

Các mệnh đề~\ref{prop:blind}--\ref{prop:coincide} cho thấy tổng lượng và độ
sâu là chưa đủ, và thứ tự thời gian chính xác là thứ đối thủ thao túng, nên
phòng thủ buộc phải quan sát dấu thời gian --- và bên duy nhất có thể bảo chứng
cho một dấu thời gian là thiết bị đã thực thi nó. Giả sử gateway giữ một bản
ghi lệnh có chữ ký
\begin{equation}
\mathcal{G}=\{(t_j,u_j,\sigma_j)\},
\label{eq:gateway}
\end{equation}
trong đó $\sigma_j$ là chữ ký trên $(t_j,u_j)$ do bộ điều khiển tạo ra tại
thời điểm đóng mở máy bơm. Phép thử neo đối chiếu nhật ký khai báo với bản ghi
đó:
\begin{equation}
\begin{aligned}
\Delta_{\mathrm{time}}&=\{t_i\}\triangle\{t_j\},\\
\Delta_{\mathrm{depth}}&=\max_{t}\bigl|u(t)-u_{\mathcal{G}}(t)\bigr|,
\end{aligned}
\label{eq:deltas}
\end{equation}
trong đó $\triangle$ là hiệu đối xứng và $\Delta_{\mathrm{depth}}$ lấy trên các
giờ có mặt ở cả hai bản ghi. Phán quyết khi đó là
\begin{equation}
A(L,\mathcal{G})=
\begin{cases}
\text{giả mạo}, & \Delta_{\mathrm{time}}\neq\emptyset
\ \vee\ \Delta_{\mathrm{depth}}>\epsilon,\\
\text{khớp}, & \text{trường hợp còn lại},
\end{cases}
\label{eq:anchor}
\end{equation}
với $\epsilon=10^{-6}$~mm, tức khớp chính xác tới mức nhiễu dấu phẩy động.
Việc gọi tên riêng hai loại sai lệch là quan trọng: dịch thời điểm chỉ làm thay
đổi $\Delta_{\mathrm{time}}$, thổi phồng độ sâu chỉ làm thay đổi
$\Delta_{\mathrm{depth}}$, còn bỏ sót và thêm mới làm thay đổi cả hai, nên một
phép đối chiếu duy nhất bắt được mọi lớp gian lận.

\begin{proposition}[Tính đầy đủ của neo thiết bị]
\label{prop:complete}
Nếu nhật ký khai báo $L$ khác bản ghi gateway $\mathcal{G}$ về đa tập hợp các
giờ hoặc về bất kỳ độ sâu nào tại một giờ chung, thì $A(L,\mathcal{G})$ trả về
\text{giả mạo}. Nếu $L$ và $\mathcal{G}$ khớp cả hai, thì $L$ là một bản sao
trung thực và không có sự giả mạo nào đối với ba đại lượng trong
\eqref{eq:three}.
\end{proposition}

\begin{proof}
Xem Phụ lục~S12: mệnh đề thứ nhất là định nghĩa của \eqref{eq:anchor}, còn
mệnh đề thứ hai suy ra được vì cả ba đại lượng có thể làm giả đều là hàm của
tập giờ và các độ sâu theo giờ.
\end{proof}

\subsection{Vì sao điều này không đưa việc kết tội oan trở lại}

Mệnh đề~\ref{prop:complete} mạnh tới mức trông nguy hiểm: mọi sai lệch đều bị
coi là giả mạo, kể cả việc quên ghi và làm tròn mà Mục~\ref{sec:guarantee} đã
chỉ ra là rất phổ biến. Cách hoá giải là ở chỗ neo thiết bị đối chiếu một bản
sao với \emph{nguồn}, chứ không đối chiếu một nhật ký với một \emph{mô hình}:
một nông dân trung thực quên một mục ghi chỉ đơn giản được bổ sung từ
$\mathcal{G}$. Khi đó, kết tội oan đòi hỏi gateway phải mâu thuẫn với chính nó
--- một vấn đề toàn vẹn thiết bị --- đó là lý do Bảng~\ref{tab:falseacc} giữ ở
mức $0/60$ trong khi các tầng hành vi giữ ở mức $26{,}7\%$ và $5{,}0\%$
(Phụ lục~S18).

\subsection{Chi phí và phạm vi của tuyên bố}

Neo thiết bị đòi hỏi gateway của công trình đồng hành~\cite{twingate2026}, vốn
đóng mở máy bơm và đã ghi một bản ghi lệnh có chữ ký kèm dấu thời gian như một
phần chức năng an toàn của nó; chi phí biên của việc chứng thực là lưu trữ và
truyền một bản ghi đằng nào cũng tồn tại (Bảng~S5 của Phụ lục). Dù vậy, phạm
vi hẹp hơn tiền đề ban đầu của chúng tôi: một nhật ký viết tay, không có thiết
bị nào trong vòng lặp, tự nó \emph{không} phải bằng chứng được chấp nhận, vì
Mục~\ref{sec:results-attack} đo được một đối thủ duy lý thổi phồng tín chỉ thêm
$+35\%$ trong khi lọt qua mọi phép thử dựa trên vật lý và tổng lượng ($6/6$
nhật ký bị sửa). MRV với chi phí gần như bằng không chỉ khả thi ở nơi có một bộ
điều khiển đứng giữa nông dân và máy bơm; ở nơi tưới thủ công, lớp này không áp
dụng được.


% BAN TIENG VIET cua 05_design_results.tex (2026-10-08). So lieu giu nguyen.
\section{Thiết kế thực nghiệm}
\label{sec:design}

Thiết kế này so sánh bốn tầng phòng thủ với chín loại nhật ký trên một thửa
ruộng cố ý vi phạm mọi giả định của mô hình sinh đôi số: khác với đánh giá
trước đây, mà $0/60$ ca kết tội oan chỉ đo tính tự nhất quán (Phụ lục~S10.4),
ruộng ở đây được sinh ra bởi một tiến trình riêng biệt, khác mô hình sinh đôi
số theo bảy trục, mỗi trục đều có nguồn đã công bố (mọi giá trị ở Bảng~S6 của
Phụ lục; danh sách từng trục và ngữ nghĩa ghi chép theo sự kiện ở S20.1). Thời
tiết là ERA5 theo giờ (kho lưu trữ Open-Meteo) cho Cần Thơ, Sóc Trăng và Cà Mau
trong năm 2024 (khan nước) và 2025 (thừa nước), cửa sổ từ ngày 1 tháng Ba đến
ngày 29 tháng Năm ($2\,160$~giờ), mùa khô mà AWD được áp dụng. Tầng~1 là bộ
phát hiện mô hình điểm đã công bố (hồ sơ mặc định, ngưỡng tuyệt đối
$\tau_0=5$~mm); Tầng~2 là bộ phát hiện bền vững của \eqref{eq:robust} với
ngưỡng tỉ lệ của \eqref{eq:tau}, chạy trên cả $\Theta$ rút gọn ($24$ lần phát
lại mỗi nhật ký, $10$ seed) và $\Theta$ đầy đủ ($150$ lần phát lại, dùng kiểm
chéo); Tầng~3 là neo thiết bị của \eqref{eq:anchor}; Tầng~4 chỉ có vật lý, dùng
làm mốc. Trong chín loại nhật ký, ba loại là giả mạo ngây thơ và bốn loại là
đối kháng, được dựng bằng một bộ tối ưu chứ không theo một công thức cố định
(danh sách đầy đủ ở S16); mọi thống kê tấn công đều có điều kiện trên cờ
\texttt{changed}, vì một đối thủ không tìm được nước đi cải thiện sẽ trả về
nguyên bản nhật ký, và các phép so sánh dùng mốc của chính biến thể đó (hai lỗi
của bộ khung thử nghiệm, S10.2). Với mỗi tầng chúng tôi báo cáo tỷ lệ lọt lưới
(nhật ký giả mà không tầng vật lý nào đánh dấu) và tỷ lệ kết tội oan trên nhật
ký trung thực; về tác động tín chỉ, báo cáo $n_{\mathrm{dry}}$ của
\eqref{eq:ndry} so với mốc ruộng trung thực, theo ghép cặp. Tổng lượng nước chỉ
xuất hiện như một chỉ số chẩn đoán (S10.3).

\section{Kết quả}
\label{sec:results}

\subsection{Bộ phát hiện đã công bố không an toàn, chứ không chỉ lạc quan}

Bảng~\ref{tab:falseacc} đưa ra kết quả phủ định trung tâm: bộ phát hiện mô hình
điểm không kết tội ai trên nhật ký điều kiện giấy, tái hiện đúng $0/60$ của lần
trước, nhưng kết tội oan $26{,}7\%$ nông dân trung thực trên nhật ký điều kiện
ruộng. Nguyên nhân là \eqref{eq:gap} với một ngưỡng tuyệt đối --- quên ghi, làm
tròn hoặc ghi trễ một ngày đều tạo ra $G>5$~mm mà không có hành vi sai trái
nào; chênh lệch đo được trên nhật ký trung thực điều kiện ruộng lên tới
$65$~mm. Ngưỡng tỉ lệ của \eqref{eq:tau} hạ con số này xuống $5{,}0\%$, và neo
thiết bị hạ xuống $0/60$.

\subsection{Vật lý không nhìn thấy tấn công sinh lợi}

Bảng~\ref{tab:credit} là kết quả phủ định thứ hai, và là kết quả quan trọng về
mặt kinh tế. Mọi con số thổi phồng đều là so sánh \emph{ghép cặp} với nhật ký
trung thực của cùng quỹ đạo, chỉ tính trên các tấn công thật sự thay đổi bản
ghi; cách tính không ghép cặp sẽ phóng đại lợi ích lên khoảng hai lần (S15).
Dịch thời điểm nâng số pha khô được cấp tín chỉ từ $4{,}78$ lên $6{,}33$ mỗi
vụ, tức $+33\%$ trên chín nhật ký bị sửa ($t=5{,}29$); với tập bất định đầy đủ,
$+38\%$ trên ba nhật ký ($t=3{,}46$); đối thủ duy lý đạt $+35\%$ và $+40\%$.
Gộp khai báo hầu như không được gì khi ghép cặp ($+1\%$, $+2\%$), và hai kiểu
giả mạo thô sơ thậm chí \emph{làm giảm} số pha khô được cấp tín chỉ ($-1\%$,
$-17\%$): chỉ có phép hoán vị dấu thời điểm giữ nguyên tổng lượng nước là sinh
lợi, đúng như Mệnh đề~\ref{prop:coincide} dự đoán. Tỷ lệ lọt lưới trong
Bảng~\ref{tab:defence} là phần cốt lõi của kết quả: dịch thời điểm lọt qua tầng
vật lý bền vững ở $8$ trong $9$ nhật ký bị sửa, đối thủ duy lý ở $6$ trong $6$.
Mô hình điểm bắt được nhiều hơn ($4$ trong $9$), nhưng nó lại chính là bộ phát
hiện kết tội oan $26{,}7\%$ nông dân trung thực; tầng bền vững hạ tỷ lệ đó xuống
$5{,}0\%$ và đổi lại giao phần lớn khả năng phát hiện về tay đối thủ ($1$ trong
$9$) --- \emph{không cấu hình vật lý nào thoát khỏi sự đánh đổi đó}, vì cả hai
tầng đều là hàm của quỹ đạo phát lại. Bộ lọc thấm $V_5$ kích hoạt $648$ lần,
nhưng chỉ trên kiểu giả mạo thổi phồng độ sâu thô sơ: hữu ích, và trực giao với
mối đe doạ sinh lợi.

\subsection{Neo thiết bị bịt mọi lớp, và cái giá của sự bền vững}

Bảng~\ref{tab:defence} cũng đưa ra nửa kiến tạo: trong $281$ nhật ký được kiểm
toán, $116$ nhật ký thật sự bị sửa, và neo thiết bị bắt được $116/116$ trong
khi không kết tội oan nhật ký nào trong $60$ nhật ký trung thực điều kiện ruộng
($37/37$ và $0/18$ với $\Theta$ đầy đủ). Giả mạo ngây thơ bị mọi tầng bắt được;
khác biệt nằm hoàn toàn ở các dòng đối kháng, nơi cả vật lý lẫn tổng lượng đều
thất bại và chỉ có neo thiết bị đứng vững. Tính bền vững không miễn phí: trên
chín nhật ký gian lận bỏ sót thật sự bị sửa, mô hình điểm bắt được $9/9$ trong
khi quy tắc bền vững bắt $4/9$ ($4/4$ so với $2/4$ ở $\Theta$ đầy đủ) --- một sự
mất mát năng lực phát hiện, không phải một cải thiện Pareto, dù chúng tôi coi
đó là phía đúng của sự đánh đổi. Mọi kết luận đều đúng với cả $\Theta$ rút gọn
($24$ lần phát lại, $281$ nhật ký, $10$ seed) và $\Theta$ đầy đủ ($150$ lần phát
lại, $85$ nhật ký, $3$ seed): kết tội oan $0/60$ và $0/18$, tính đầy đủ của neo
$116/116$ và $37/37$, lọt lưới dịch thời điểm $8/9$ và $3/3$, thổi phồng ghép
cặp $+33\%$ và $+38\%$ (Phụ lục~S17).


\begin{table}[!t]
\centering
\caption{Kết tội oan trên nhật ký trung thực, thuộc tính mà cả phương án được
chào bán. Điều kiện giấy tái hiện kết quả $0/60$ trước đây; điều kiện ruộng
phơi bày ngưỡng chênh lệch tuyệt đối.}
\label{tab:falseacc}
\footnotesize
% Sinh TU DONG tu tables/tab_false_accusation.tex (ban tieng Anh, von sinh tu outputs/*.csv)
% boi experiments/make_tables_hujos_vi.py. KHONG sua tay.
% Nhan dich qua bang tra cuong che; so chi doi dinh dang sang kieu Viet Nam (dau phay thap phan).
% Da assert: cung so dong, cung so vach, cung tap chu so voi ban EN.

% Sinh TU DONG tu outputs/eval_theta_*.csv boi experiments/make_tables_anchored.py
% KHONG sua tay: moi con so deu truy nguon duoc ve CSV.

\fittocol{%
\begin{tabular}{@{}>{\raggedright\arraybackslash}p{4.20cm}>{\raggedright\arraybackslash}p{2.10cm}>{\raggedright\arraybackslash}p{2.10cm}@{}}
\toprule
Lớp phòng thủ & Trung thực, điều kiện giấy & Trung thực, điều kiện ruộng \\
\midrule
v1 tầng 0 (vật lý, mô hình điểm) & 0/60 & 2/60 \\
v1 tầng 1 (tuyệt đối $G>5$\,mm) & 0/60 & 16/60 \\
v2 tầng 0 ($\forall$ trên $\Theta$) & 0/60 & 0/60 \\
v2 tầng 1 ($\tau(n)$ co giãn theo số lệnh) & 0/60 & 3/60 \\
Bản ghi lệnh neo theo thiết bị & 0/60 & 0/60 \\
\bottomrule
\end{tabular}
}

\end{table}

\begin{table}[!t]
\centering
\caption{Số pha khô được cấp tín chỉ dài ít nhất $72$~giờ --- đại lượng mà
VM0051 và Quyết định 4801 trả tiền --- so với mốc ruộng trung thực (ghép cặp).
Cột cuối đếm số tấn công thật sự sửa nhật ký. Bảng có hai mốc: dòng ``trung
thực'' là trung bình $3{,}23$ pha của toàn bộ $60$ nhật ký trung thực, còn
cột ``trước'' của mỗi dòng gian lận là nhật ký trung thực \emph{ghép cặp} cùng
trạm, năm, seed và loại đất ($4{,}87$ với các biến thể đối kháng) --- hai mốc
khác nhau có chủ đích, xem Phụ lục~S15.}
\label{tab:credit}
\footnotesize
% Sinh TU DONG tu tables/tab_credit.tex (ban tieng Anh, von sinh tu outputs/*.csv)
% boi experiments/make_tables_hujos_vi.py. KHONG sua tay.
% Nhan dich qua bang tra cuong che; so chi doi dinh dang sang kieu Viet Nam (dau phay thap phan).
% Da assert: cung so dong, cung so vach, cung tap chu so voi ban EN.

% Sinh TU DONG tu outputs/eval_theta_*.csv boi experiments/make_tables_anchored.py
% KHONG sua tay: moi con so deu truy nguon duoc ve CSV.

\fittocol{%
\begin{tabular}{@{}>{\raggedright\arraybackslash}p{2.90cm}c>{\raggedright\arraybackslash}p{1.90cm}cc@{}}
\toprule
Loại nhật ký & Nhật ký bị sửa & Pha khô ghép cặp (trước \to sau) & Thổi phồng & $t$ \\
\midrule
Nông dân trung thực (điều kiện ruộng) & --- & 3,23 (mốc) & mốc & --- \\
Thêm lệnh tưới & 23/23 & 4,87 \to 4,83 & -1\% & -0,57 \\
Độ sâu khai không khả thi & 23/23 & 4,87 \to 4,04 & -17\% & -4,03 \\
Bỏ sót lệnh thật & 9/23 & 5,44 \to 4,78 & -12\% & -2,83 \\
Gộp khai báo, 24 h & 23/23 & 4,87 \to 4,91 & +1\% & 0,37 \\
Gộp khai báo, 96 h & 23/23 & 4,87 \to 4,96 & +2\% & 0,62 \\
Dịch thời điểm (duy lý) & 9/23 & 4,78 \to 6,33 & +33\% & 5,29 \\
Đối thủ duy lý & 6/23 & 4,33 \to 5,83 & +35\% & 4,39 \\
\bottomrule
\end{tabular}
}

\end{table}

\begin{table}[!t]
\centering
\caption{Kết quả chính. Mọi tỷ lệ được tính trên các nhật ký mà tấn công
\emph{thật sự} sửa; cột thứ hai báo số bị sửa trên số được sinh, vì đối thủ trả
về nguyên bản nhật ký khi không tìm được nước đi cải thiện. Tỷ lệ lọt lưới là
phần mà không tầng vật lý nào đánh dấu.}
\label{tab:defence}
\footnotesize
% Sinh TU DONG tu tables/tab_defence.tex (ban tieng Anh, von sinh tu outputs/*.csv)
% boi experiments/make_tables_hujos_vi.py. KHONG sua tay.
% Nhan dich qua bang tra cuong che; so chi doi dinh dang sang kieu Viet Nam (dau phay thap phan).
% Da assert: cung so dong, cung so vach, cung tap chu so voi ban EN.

% Sinh TU DONG tu outputs/eval_theta_*.csv boi experiments/make_tables_anchored.py
% KHONG sua tay: moi con so deu truy nguon duoc ve CSV.

\fittocol{%
\begin{tabular}{@{}>{\raggedright\arraybackslash}p{2.85cm}c>{\raggedright\arraybackslash}p{1.30cm}>{\raggedright\arraybackslash}p{1.50cm}>{\raggedright\arraybackslash}p{1.45cm}@{}}
\toprule
Biến thể gian lận & Nhật ký bị sửa & Lọt v1 & Lọt v2 vật lý & Neo thiết bị bắt \\
\midrule
Nông dân trung thực (điều kiện ruộng) & 0/60 (trung thực) & \textbf{5\%} bị oan & 0\% & 0\% bị oan \\
Thêm lệnh tưới & 23/23 & 0\% & 0\% & 23/23 \\
Độ sâu khai không khả thi & 23/23 & 0\% & 0\% & 23/23 \\
Bỏ sót lệnh thật & 9/23 & 0\% & 56\% & 9/9 \\
Gộp khai báo, 24 h & 23/23 & 39\% & 91\% & 23/23 \\
Gộp khai báo, 96 h & 23/23 & 39\% & 91\% & 23/23 \\
Dịch thời điểm (duy lý) & 9/23 & 56\% & 89\% & 9/9 \\
Đối thủ duy lý & 6/23 & 100\% & 100\% & 6/6 \\
\midrule
\textbf{Mọi nhật ký thật sự giả} & \textbf{116} & \textbf{---} & \textbf{---} & \textbf{116/116} \\
\bottomrule
\end{tabular}
}

\end{table}

% Bang benchmark: so sanh voi cac phuong an MRV da cong bo (SOTA).
% So sinh tu experiments/benchmark_confirmation.py -> _benchmark_summary.json
% KHONG sua tay; nhan dich boi make_tables_hujos_vi.py tu bang EN.
\begin{table}[!t]
\centering
\caption{Số pha khô đủ điều kiện cấp tín chỉ có thể xác nhận được theo từng
phương án MRV, trên $150$ pha có thật trong $60$ ruộng mô phỏng; một pha chỉ
được tính nếu có ít nhất một lần quan sát rơi vào trong nó. Hai dòng cuối hoà
nhau ở chỉ số này và khác nhau về chi phí cũng như về thứ chúng chứng minh
được với một nhật ký đã bị sửa.}
\label{tab:benchmark}
\footnotesize
% Sinh TU DONG tu tables/tab_benchmark.tex (ban tieng Anh, von sinh tu outputs/*.csv)
% boi experiments/make_tables_hujos_vi.py. KHONG sua tay.
% Nhan dich qua bang tra cuong che; so chi doi dinh dang sang kieu Viet Nam (dau phay thap phan).
% Da assert: cung so dong, cung so vach, cung tap chu so voi ban EN.

% Sinh TU DONG tu outputs/_benchmark_summary.json boi experiments/make_table_benchmark.py
% KHONG sua tay: moi con so truy nguon ve benchmark_confirmation.py.

\fittocol{
\begin{tabular}{@{}>{\raggedright\arraybackslash}p{4.30cm}ccc@{}}
\toprule
Phương án xác minh & Pha khô xác nhận được ($n=150$) & Số lần quan sát mỗi vụ & Thiết bị mỗi thửa \\
\midrule
Ảnh + tự khai (thực hành hiện hành)~\cite{qd4801} & 0/150 (0\%) & --- & 0 (ống PVC) \\
Sentinel-1A+1B (chu kỳ 6 ngày)~\cite{hashemi2022s1} & 124/150 (83\%) & 15 & 0 \\
Sentinel-1A đơn lẻ (chu kỳ 12 ngày)~\cite{clauss2017s1} & 90/150 (60\%) & 8 & 0 \\
Ống đo mực nước liên tục (IoT)~\cite{delacruz2022} & 150/150 (100\%) & 2,153 & 1 mỗi thửa \\
Bản ghi lệnh có chữ ký (bài này)~\cite{twingate2026} & 150/150 (100\%) & 2,153 & 0 (bộ điều khiển) \\
\bottomrule
\end{tabular}
}

\end{table}

% BAN TIENG VIET cua 06_discussion_conclusion.tex (2026-10-08). So lieu giu nguyen.
\section{Bàn luận}
\label{sec:discussion}

\subsection{Quan hệ với các công trình trước}

Việc giám sát AWD đã được tiếp cận từ ba hướng --- cảm biến mực nước có thiết
bị~\cite{delacruz2022}, viễn thám~\cite{lovell2019,islam2026alose,mlrice2025},
và MRV ở cấp chương trình~\cite{qd1490,mae2026progress,irri2026mrv} --- trong
khi phân tích tổng hợp nông học~\cite{carrijo2017} và các nghiên cứu về tiếp
nhận~\cite{lampayan2015,enriquez2021} giải thích vì sao giám sát có thiết bị
không lan rộng được. Không hướng nào đánh giá \emph{độ tin cậy của chính bản
ghi hành động}, và đó là câu hỏi được đặt ra ở đây (khảo sát đầy đủ ở
Phụ lục~S8). Với radar vệ tinh, phương án chi phí thấp hiển nhiên, chúng tôi đo
chứ không tranh luận: Bảng~\ref{tab:benchmark} chấm điểm mọi phương án theo số
pha khô đủ điều kiện cấp tín chỉ mà nó xác nhận được trong $150$ pha, một pha
chỉ được tính nếu có ít nhất một lần quan sát rơi vào trong
nó~\cite{clauss2017s1,hashemi2022s1,delacruz2022}. Hai điểm cần thận trọng: cột
này đo khả năng \emph{kiếm} tín chỉ chứ không đo khả năng phát hiện gian lận ---
ống đo và bản ghi lệnh cùng đạt $100\%$ nhưng khác nhau về chi phí và về thứ
chúng chứng minh được --- và ảnh chụp kèm tự khai đạt $0\%$ vì nó không tạo ra
chuỗi quan sát nào để đối chiếu. Do đó radar vẫn giữ vai trò mà Phụ lục 4 của
VM0051 giao cho nó: một phép kiểm chéo theo thửa làm giảm tần suất kiểm tra thực
địa~\cite{vm0051} (Phụ lục~S9.2).

\subsection{Giới hạn}

Bảy giới hạn ràng buộc các tuyên bố, tất cả đều được phát biểu đầy đủ ở
Phụ lục~S21; hai giới hạn chi phối các kết quả chính là như sau. \emph{Nông dân
mô phỏng}: các bản giả mạo được tiêm vào quỹ đạo mô phỏng, nên vẫn cần một thử
nghiệm thực địa để đo phản ứng hành vi --- trên hết là liệu nông dân có thay đổi
việc bơm \emph{thật sự} của mình một khi biết có neo thiết bị hay không, một tấn
công làm thay đổi $M$ chứ không phải $L$ và mô hình của chúng tôi không biểu
diễn được. \emph{Tấn công chỉ được minh hoạ trong năm khan nước}: năm 2025 phần
lớn nhật ký rơi xuống dưới mức tối thiểu ba lệnh mà tại đó một tấn công thời điểm
mới có thể dựng được --- trong $37$ nhật ký thật sự bị sửa ở $\Theta$ đầy đủ, cả
$37$ đều thuộc năm 2024. Chính vách ngăn giữa hai năm đó là một phát hiện (rủi
ro gian lận MRV tập trung ở nơi khan nước), nhưng kết quả tấn công chưa phải là
bằng chứng về một mùa mưa; ở $\Theta$ rút gọn có mười hai nhật ký bị sửa thuộc
năm 2025 và chúng hành xử giống nhau --- gợi ý chứ chưa kết luận. Năm giới hạn
còn lại được liệt kê ở S21.

\section{Kết luận}

Chúng tôi khởi đầu với mục tiêu làm cho MRV lúa gần như miễn phí bằng cách kiểm
toán một nhật ký tưới do nông dân tự ghi đối chiếu với vật lý nước trong đất.
Phiên bản đầu tiên của lớp đó báo cáo không có kết tội oan và phát hiện trọn
vẹn; cả hai con số đều là sản phẩm của một thửa ruộng do chính mô hình kiểm
toán sinh ra, và việc dựng lại thửa ruộng đó trước một đối thủ tối ưu hoá đại
lượng được cấp tín chỉ đã làm thay đổi kết luận.

Bộ phát hiện mô hình điểm kết tội oan $26{,}7\%$ nông dân trung thực; một ngưỡng
tỉ lệ và một phép giao trên tập bất định $150$ phần tử hạ con số đó xuống
$5{,}0\%$ và $0/60$, với cái giá đo được về năng lực phát hiện ($9/9$ xuống
$4/9$ ở gian lận bỏ sót). Trong khi đó một phép hoán vị thời điểm khai báo giữ
nguyên tổng lượng nước làm thổi phồng số pha khô được cấp tín chỉ từ $33\%$ đến
$38\%$, đồng thời lọt qua tầng bền vững ở $8$ trong $9$ và $3$ trong $3$ nhật ký
bị sửa, và không cấu hình vật lý nào thoát khỏi sự đánh đổi đó: hoãn một khai báo
vừa kéo dài pha khô vừa làm tăng phần thiếu hụt dành cho nó, nên lợi nhuận và sự
không thể phát hiện dùng chung một gradient.

Cách sửa không thuộc về thuỷ văn. Đối chiếu nhật ký khai báo với bản ghi lệnh có
chữ ký và dấu thời gian của gateway bắt được $116/116$ nhật ký thật sự bị giả
mạo và không kết tội oan nhật ký nào trong $60$ nhật ký trung thực, với chi phí
không vượt quá một bộ điều khiển mà công trình đồng hành đã triển khai. Điều còn
lại của tiền đề ban đầu là: vật lý \emph{ràng buộc} một nhật ký, nhưng chỉ một
bản ghi thiết bị có chữ ký mới \emph{chứng thực} được nó --- MRV chi phí gần như
bằng không chỉ khả thi ở nơi có một thiết bị đứng giữa nông dân và máy bơm. Ba
con số đã công bố trước đó được sửa lại tương ứng ($[49{,}0; 86{,}0]$~mm thay vì
$61$~mm; một bộ lọc thứ ba bất động; một phép thử chênh lệch tuyệt đối không thể
chấp nhận). Mã nguồn, cả hai tập bất định, toàn bộ $281$ nhật ký đã kiểm toán và
$46$ phép thử đơn vị đều được công khai.


% ---------------------------------------------------------------- khai bao
\section*{Dữ liệu, mã nguồn và tuyên bố}

Mã nguồn, cả hai tập bất định, cả hai cài đặt đối thủ và toàn bộ $281$ nhật ký
đã kiểm toán có tại \url{https://github.com/mxuanvan02/TwinGate} (giấy phép
MIT); dữ liệu khí tượng là tái phân tích ERA5 công khai qua kho lưu trữ
Open-Meteo. Nghiên cứu này được tài trợ bởi Đại học Huế trong đề tài mã số
DHH2025-19-07. Các tác giả tuyên bố không có xung đột lợi ích. Không áp dụng
phê duyệt đạo đức: nghiên cứu chỉ dùng dữ liệu tái phân tích công khai và thí
nghiệm mô phỏng, không có người tham gia, động vật hay can thiệp thực địa.

% ---------------------------------------------------------------- tai lieu
\bibliographystyle{vancouver}
\bibliography{refs}

% ---------------------------------------------------------------- thong tin tac gia
\vspace{10pt}
\begingroup
\normalfont\fontsize{9}{11}\selectfont
\noindent\textbf{Thông tin tác giả.}\\[2pt]
\noindent\textbf{Mai Xuân Văn.} Cử nhân. Trường THPT Thuận Hóa, Trường Đại học
Sư phạm, Đại học Huế. Số 34 Lê Lợi, phường Thuận Hóa, thành phố Huế, Việt Nam.\\[2pt]
\noindent\textbf{Nguyễn Tương Tri.} Tiến sĩ. Tác giả liên hệ. Trường Đại học Sư
phạm, Đại học Huế; Viện Đào tạo Mở và Công nghệ Thông tin, Đại học Huế. Số 34
Lê Lợi, phường Thuận Hóa, thành phố Huế, Việt Nam. Email: ntuongtri@hueuni.edu.vn.\\[2pt]
\noindent Nghiên cứu này được tài trợ bởi Đại học Huế trong đề tài mã số
DHH2025-19-07.\\[4pt]
\noindent\textit{Chữ ký tác giả:} \rule{4cm}{0.4pt}\hfill\rule{4cm}{0.4pt}
\par
\endgroup

% ---------------------------------------------------------------- khoi tieng Anh (cuoi bai)
\vspace{10pt}
\begingroup
  \hujostitle{Device-Anchored Irrigation Logs: An Adversarial Evaluation of
  Near-Zero-Cost MRV for Low-Emission Rice}

  \vspace{8pt}
  \hujosauthors{Xuan Van Mai$^{1}$, Tuong Tri Nguyen$^{2,3,*}$}

  \vspace{4pt}
  \hujosaffil{$^{1}$ Thuan Hoa High School, University of Education, Hue University, Hue, Viet Nam\\
  $^{2}$ Institute of Open Education and Information Technology, Hue University, Hue, Viet Nam\\
  $^{3}$ University of Education, Hue University, 34 Le Loi St., Hue, Viet Nam}

  \vspace{4pt}
  \hujosaffil{$^{*}$ Corresponding author: Tuong Tri Nguyen. E-mail: ntuongtri@hueuni.edu.vn}

  \vspace{6pt}
  \begin{hujosabstract}
  \noindent\textbf{Abstract.}

  Monitoring, reporting and verification (MRV) is the bottleneck keeping emission
    reductions in low-emission rice from becoming tradable carbon credits, and
    monitoring is about three quarters of that cost. Because every credit-eligible
    action under alternate wetting and drying is already an \emph{actuation command},
    a farmer-written irrigation log looks like free evidence of the highest rank under
    Decision 4801/QD-BNNMT. We built that layer on an FAO-56 digital twin, then
    attacked it; the attack is the contribution. Replaying a log under public ERA5
    forcing exposes claimed depths no soil can hold, but an adversary who keeps the
    total volume and every individual depth and merely \emph{shifts commands in time}
    inflates creditable dry phases by $33\%$ in paired comparison. No physical tier
    escapes it usefully: the point model accuses $26.7\%$ of honest farmers, and the
    robust tier that cuts false accusation to $5.0\%$ in turn misses most
    time-shifting attacks. Over $281$ audited logs at three Mekong Delta stations
    and two climate years, matching the log against the gateway's signed,
    time-stamped command record catches $116/116$ forged logs and accuses none.
    Physics-based attestation cannot stand alone; the timestamp must come from
    a device, not the farmer.

\noindent\textbf{Keywords:} alternate wetting and drying; monitoring, reporting
  and verification; carbon credit; low-emission rice; digital twin.

  \vspace{10pt}
  \end{hujosabstract}
\endgroup

\end{document}
